\section{Discussion and Conclusion}
\label{sec:conclusion}

% -----------------------------------------------------------------------------
\subsection{Discussion}

\subsubsection{Clinical Significance of Over-Segmentation Suppression}
In clinical computer-assisted diagnosis, high sensitivity is a fundamental baseline prerequisite, but low specificity and rampant over-segmentation represent the primary barriers preventing clinical adoption~\cite{leufkens2012factors, mori2019real}. False-positive alarms---where healthy mucosa, vascular reflections, or background tissue are flagged as malignant lesions---induce severe cognitive fatigue in physicians and radiologists~\cite{mori2019real, hassan2021overcoming}. Crucially, over-segmented lesion boundaries mislead clinicians during surgical excision or endoscopic resection, elevating procedural risks such as tissue perforation and hemorrhage~\cite{rex2017colorectal}.

Our proposed GDF-Net architecture directly addresses this clinical hazard across multiple imaging modalities. By establishing the Feedback Firewall, GDF-Net eliminates the echo chamber of recurrent error propagation. On the CHASE-DB1 retinal benchmark, GDF-Net cuts the false-positive rate by \textbf{65\%} (reducing FPR from 2.14\% to 0.74\%) while elevating precision from 81.70\% to \textbf{86.37\%}. On CVC-ClinicDB colonoscopy polyps, FPR is reduced by over $2.4\times$ (down to 0.40\%), producing tightly bounded, clinically dependable anatomical segmentations.

\subsubsection{Universal Architectural Principle: Decoupling Forward Guidance from Backward Optimization}
Our theoretical analysis reveals a fundamental, universal design principle for recurrent vision architectures: \textbf{forward iterative spatial guidance must be decoupled from backward gradient propagation}. In conventional recurrent feedback designs, backpropagating through recurrent predictions creates a closed, non-stationary loop: the network updates its internal feature extractors to validate its own past predictions rather than the underlying anatomy.

By inserting the stop-gradient operator $\text{detach}(m^{(t-1)})$, we transform historical predictions into clean, uncorrupted spatial conditioning priors. Concurrently, our Multi-scale Gradient-Decoupled Feedback Decoder ensures dense, continuous gradient flow across all hierarchical decoder resolutions ($\mathbf{F}_\ell \odot (1 + \mathbf{S}_\ell)$). This combination preserves feature stability throughout early encoder stages and prevents the catastrophic representation drift that plagued earlier feedback networks.

\subsubsection{Addressing Low-Data Regimes via Transfer Learning}
A key observation from our benchmark across the 7 datasets is the vulnerability of recurrent feedback architectures in extreme low-data regimes. On the DRIVE retinal vessel cohort (20 training images), training a recurrent feedback network from scratch leads to catastrophic representational collapse (0.2837 Dice). In contrast, pre-training the encoder backbone on a large medical cohort (ISIC-2018 melanoma) and transferring weights to DRIVE unlocks rapid convergence to \textbf{0.7932} Dice, surpassing the published baseline across all metrics. This proves that while recurrent attention requires sufficient feature diversity to avoid overfitting, transfer learning effectively bridges the sample-efficiency gap in low-data biomedical domains.

\subsubsection{Limitations and Future Work}
While GDF-Net demonstrates consistent superiority across all seven 2D biomedical benchmarks, several promising avenues remain for future exploration:
\begin{enumerate}
  \item \textbf{Extension to 3D Volumetric Modalities:} Extending GDF-Net and the Feedback Firewall to 3D medical imaging volumes (e.g., multi-contrast brain MRI, abdominal CT) where recurrent slice-to-slice feedback could provide inter-slice volumetric continuity.
  \item \textbf{Real-Time Video Colonoscopy Tracking:} Integrating temporal recurrent gating with optical flow estimation for continuous video polyp tracking, enabling real-time frame rates ($>30\text{ FPS}$) during live colonoscopy procedures.
  \item \textbf{Adaptive Multi-Scale Gating:} Investigating dynamic early-stopping mechanisms based on prediction entropy between feedback steps $t=1$ and $t=2$ to accelerate inference on unambiguous anatomical targets.
\end{enumerate}

% -----------------------------------------------------------------------------
\subsection{Conclusion}
In this paper, we identified, diagnosed, and resolved the \textbf{Feedback Trap}---a pervasive structural pathology in recurrent medical vision networks where prediction history acting simultaneously as forward spatial guidance and backward optimization causes catastrophic error amplification and severe false-positive over-segmentation.

To overcome this dilemma, we formulated the \textbf{Feedback Firewall} principle and introduced \textbf{GDF-Net}, an architecture integrating a 4-channel decoupled guidance stem with a Multi-Scale Gradient-Decoupled Feedback Decoder. Extensive empirical benchmarks across seven datasets spanning four biomedical modalities (colonoscopy polyps, cell membranes, skin lesions, retinal vessels, and nuclei) prove that GDF-Net outperforms standard baseline architectures across 100\% of evaluated benchmarks, achieving \textbf{0.9427} Dice on CVC-ClinicDB (+4.69\,pp), \textbf{0.9518} on EM-Dataset (+17.02\,pp), \textbf{0.8880} on Kvasir-SEG (+0.28\,pp), \textbf{0.9000} on ISIC-2018 (+2.69\,pp), \textbf{0.8953} on DSB-2018 (+2.32\,pp), and slashing false-positive rates by \textbf{65\%} on CHASE-DB1. Furthermore, our theoretical analysis and empirical ablations prove that gradient decoupling is the fundamental prerequisite for robust, stable recurrent learning in biomedical computer vision.
